\documentclass{ximera}
\graphicspath{ {./images/} }
\usepackage{caption}

\begin{document}
\captionsetup{singlelinecheck=false}
Print Name: $\_\_\_\_$\\
OSU Name.\#: $\_\_\_\_$\\
Instructor: $\_\_\_\_$\\
Signature: $\_\_\_\_$

\section*{MATH 1075}
\section*{Midterm Exam 3}
\section*{Autumn 2019}
\section*{Instructions:}
\begin{itemize}
  \item You have 55 minutes to complete this exam.
  \item When applicable, mark answer bubbles completely like this , NOT like this or .
  \item Circle or box all other final answers.
  \item A calculator may be used given the calculator policy outlined in the syllabus. Even if a calculator is used, you must show all work on each problem to receive full credit.
  \item Write clearly and legibly to receive credit.
  \item Do not round; give only exact answers (unless specified otherwise).
\end{itemize}

\section*{Exam 3 - Form A}
Problem 1. Graphs of two functions, $f(x)$ and $g(x)$, are given below. The graph of $f(x)$ is solid, and the graph of $g(x)$ is dashed. Use the graphs of the two functions to fill in the blanks below.\\
(12 pts)\\
\includegraphics[alt={Graph showing two piecewise linear functions y equals f(x) as a solid line and y equals g(x) as a dashed line},width=\textwidth]{piecewise-functions-f-and-g.jpg}

$$\begin{aligned}
& g(0)= \\
& (f+g)(-1)= \\
& (f-g)(1.5)=
\end{aligned}$$

$$\begin{aligned}
& (f g)(3)= \\
& (g \div f)(0)= \\
& (g \circ f)(0)=
\end{aligned}$$

Problem 2. Graphs of two relations, Relation A and Relation B, are given below. (8 pts)\\
\includegraphics[alt={Graph of Relation A showing six plotted points on a coordinate plane},width=\textwidth]{relation-A-graph}

\begin{figure}[h]
\begin{center}
  \includegraphics[alt={Graph of Relation B showing a piecewise function with open and closed circles indicating excluded and included points},width=\textwidth]{relation-B-graph}
\captionsetup{labelformat=empty}
\caption{Relation B}
\end{center}
\end{figure}

(a) Select the statement that is true.\\
◯ Relation A is a function, but Relation B is not.\\
◯ Relation B is a function, but Relation A is not.\\
◯ Both Relation A and Relation B are functions.\\
◯ Neither Relation A nor Relation B are functions.\\
(b) Which of the following represents the domain of Relation A?\\
◯ $[1,2]$\\
◯ $\{-1,1,2\}$\\
◯ $[1,4]$\\
◯ $\{0,1,2,3,4\}$\\
(c) Which of the following represents the range of Relation B?\\
◯ $\{1,2,3\}$\\
◯ $[1,3)$\\
◯ $\{1,2,3,4\}$\\
◯ $(1,4)$\\
(d) Fill in the blank so that the ordered pair $(x, y)$ is included in Relation B: ( $\_\_\_\_$ 2). Select your answer from the options below.\\
◯ 1\\
◯ 2\\
◯ 3\\
◯ 4

\section*{Exam 3 - Form A}
Problem 3. Simplify. Assume that the variable represents a positive real number. (8 pts)

$$\left(\frac{1}{32 x^{5}}\right)^{-2 / 5}$$

Problem 4. Simplify. Assume all variables represent positive real numbers. (10 pts)

$$\sqrt[4]{81 x^{19} z^{17}}$$

Problem 5. Simplify. Assume all variables represent positive real numbers. (8 pts)

$$x \sqrt{75 x v^{2}}+7 v \sqrt{27 x^{3}}$$

\section*{Exam 3 - Form A}
Problem 6. Solve for $y$, where $y$ is a real number.\\
(11 pts)

$$y-4=\sqrt{-4 y+37}$$

Problem 7. Solve for $w$, where $w$ is a real number.

$$(w-7)^{2}-20=0$$

\section*{Exam 3 - Form A}
Problem 8. Simplify. Assume the variable represents a positive real number. (10 pts)

$$2 \sqrt[3]{500 x^{4}}-\sqrt[3]{256 x^{4}}$$

Problem 9. Simplify. Assume all variables represent positive real numbers. Be sure to rationalize the denominator.

$$\sqrt{\frac{7 h^{2} k^{4}}{32}}$$

\section*{Exam 3 - Form A}
Problem 10. Simplify. Assume the variable represents a positive real number. (8 pts)\\
$10 \sqrt[5]{x^{4}} \cdot 6 \sqrt{x}$

Problem 11. Rationalize the denominator and simplify. Assume that the variable represents a positive real number.\\
(8 pts)\\
$\frac{-3}{3 \sqrt{w}+4}$


\end{document}